\documentclass[11pt]{article}
\usepackage{../shared/luxcover}
\usepackage{amsmath,amssymb}
\usepackage{hyperref}
\usepackage{booktabs}
\usepackage{enumitem}
\usepackage{xcolor}

\title{Lux Quantum Custody}
\author{Zach Kelling\\\textit{Lux Industries, Inc.}\\\texttt{research@lux.network}}
\date{January 2025}

\begin{document}

\luxcoverpage

\begin{abstract}
This paper presents the design and architecture of the described system.
\end{abstract}

\tableofcontents
\newpage


\subsection{The Value Proposition of Lux Network}

\subsection{Why should I wrap and bridge my Bitcoin onto the Lux network?}

Opting for this approach presents a strategic choice that provides heightened protection against quantum computing risks, allows you to stake your BTC to earn yield.


BTC holders on the Lux Network can utilize stablecoins obtained through DeFi loans to acquire gold and silver Digital-Safe Keeping Receipts (D-SKR's), thus diversifying their investment portfolio and creating a bridge between digital assets and tangible real-world commodities.


\subsection{Enhanced Security with Quantum Resistance}

As a BTC holder, moving your assets to the Lux Network offers a significant upgrade in security, especially in the context of emerging quantum computing threats.


\subsubsection{SHA-256 vs. Lamport Signatures}

SHA-256 in Bitcoin: Bitcoin uses the SHA-256 hashing algorithm, which is fundamental to its mining and transaction validation processes. While SHA-256 is robust against current computational capabilities, its long-term resilience against quantum computing is uncertain.


Lamport Signatures in Lux Network: Lamport signatures, as part of the Lux Network's quantum-resistant approach, offer a higher degree of security against quantum computing threats. They are based on simple mathematical operations and are considered one of the strong candidates for post-quantum cryptography.


\subsubsection{Lattice-Based MPC and Homomorphic Encryption}

Lattice-Based MPC: Multi-Party Computation (MPC) based on lattice cryptography in the Lux Network allows for secure, collaborative computation among multiple parties without revealing individual inputs. This is crucial for maintaining privacy and security in decentralized finance transactions and smart contract executions.


Homomorphic Encryption: This form of encryption allows computations to be carried out on encrypted data, providing results that, when decrypted, match the outcome of operations performed on the plaintext. In the Lux Network, this means enhanced data privacy and security for transactions and smart contracts.


\subsection{Leveraging BTC on Lux Network}

Non-Liquidation Asset Leverage and Borrowing: By bridging BTC to the Lux Network, you can stake your BTC as collateral to borrow Lux coins. This mechanism allows you to maintain your Bitcoin investment while accessing additional liquidity.


High-Yield Staking Opportunities: Staking the borrowed Lux coins on a Lux validator can yield substantial returns, notably a 11\% return over a 12-month staking period.


\subsection{Implications for BTC Holders}

Future-Proofing Investments: Transitioning from Bitcoin's SHA-256 to the quantum-resistant technologies like Lamport signatures, lattice-based MPC, and homomorphic encryption on the Lux Network is a forward-looking step, ensuring the long-term security of your assets.


Expanding Investment Horizons with D-SKR's: The opportunity to mint asset-backed digital currencies using D-SKR's expands your investment horizons, allowing you to link your digital holdings with physical assets, adding a layer of stability and tangible value to your portfolio.


\subsection{Conclusion}

In conclusion, bridging BTC to the Lux Network is a strategic decision that offers enhanced security against quantum threats, a diversified investment portfolio through D-SKR's, and lucrative staking opportunities. It represents a significant evolution from the SHA-256-based security model of Bitcoin, embracing advanced cryptographic techniques that are poised to withstand the advent of quantum computing, while also expanding the utility and potential of your digital assets.


\end{document}
